% Figura 2.1 — La imagen muográfica por transmisión como proyección (esquema sin escala).
% Fuente editable (TikZ). Compilar con:  pdflatex fig2-1_opacidad_proyeccion.tex
% (a) Cada dirección de observación (acimut alpha, elevación eta) integra la densidad a lo
%     largo de un trayecto: X = integral de rho(l) dl. Un píxel resume todo el trayecto.
% Notación igual a la del capítulo 3 (alpha, eta, X, rho, l, L).
% (b) Ambigüedad espesor-densidad: tres columnas con la misma opacidad
%     X = 1,06e5 g/cm2: 400 m de roca de 2,65 g/cm3; 530 m de 2,00 g/cm3;
%     450 m con una cavidad de 50 m y 400 m de roca de 2,65 g/cm3.
% Paleta Paul Tol "bright", igual que las figuras del capítulo 3.
\documentclass[12pt,border=2pt]{standalone}
\usepackage[T1]{fontenc}
\usepackage[utf8]{inputenc}
\usepackage[spanish]{babel}
\usepackage{lmodern}
\usepackage{amsmath}
\usepackage{tikz}
\usetikzlibrary{babel,arrows.meta,calc,intersections,decorations.pathreplacing}

\definecolor{tolblue}{HTML}{4477AA}
\definecolor{tolgreen}{HTML}{228833}
\definecolor{tolyellow}{HTML}{CCBB44}
\definecolor{tolred}{HTML}{EE6677}
\definecolor{tolpurple}{HTML}{AA3377}
\definecolor{tolgrey}{HTML}{BBBBBB}

\begin{document}
\footnotesize
\begin{tikzpicture}[>={Stealth[length=2mm]}, etiqueta/.style={font=\scriptsize, text=black, align=center, inner sep=1pt}]

% =========================== (a) Trayectos y proyección ===========================
\begin{scope}[scale=0.72]
  \coordinate (det) at (5.4,0.12);
  % Terreno y edificio
  \draw[black!45, line width=0.5pt] (-3.3,0) -- (6.3,0);
  \filldraw[name path=cono, fill=tolgrey!40, draw=black!70, line width=0.6pt] (-3,0) -- (0,2.6) -- (3,0) -- cycle;
  % Conducto de baja densidad y cuerpo denso
  \filldraw[fill=tolblue!45, draw=tolblue!80!black, line width=0.4pt] (0.3,1.0) circle (0.32);
  \filldraw[fill=tolred!55, draw=tolred!80!black, line width=0.4pt] (-0.85,0.48) ellipse (0.42 and 0.24);
  \node[etiqueta, text=tolblue!80!black, anchor=west] at (0.72,0.62) {conducto};
  \node[etiqueta, text=tolred!70!black] at (-2.3,-0.32) {cuerpo denso};
  \draw[tolred!70!black, line width=0.3pt] (-1.75,-0.22) -- (-1.1,0.32);

  % Trayecto por cielo abierto (referencia)
  \draw[tolgreen!80!black, line width=0.5pt, dash pattern=on 2pt off 1.5pt] (-1.0,4.2) -- (det);
  \node[etiqueta, text=tolgreen!60!black, anchor=south] at (-1.0,4.2) {cielo abierto};

  % Trayecto inferior, a través del cuerpo denso
  \draw[black!55, line width=0.4pt] (-3.3,0.536) -- (det);

  % Trayecto principal, a través del conducto
  \path[name path=rayo] (det) -- (-3.3,1.86);
  \path[name intersections={of=cono and rayo, sort by=rayo}]
    (intersection-1) coordinate (entra)
    (intersection-2) coordinate (sale);
  \draw[tolblue!85!black, line width=0.4pt, dash pattern=on 1.5pt off 1.5pt] (-3.3,1.86) -- (sale);
  \draw[tolblue!85!black, line width=1.0pt] (sale) -- (entra);
  \draw[tolblue!85!black, line width=0.4pt, dash pattern=on 1.5pt off 1.5pt] (entra) -- (det);
  \fill[tolblue!85!black] (entra) circle (1.3pt) (sale) circle (1.3pt);
  % Longitud del recorrido en la roca
  \draw[decorate, decoration={brace, amplitude=4pt, raise=3pt}, black!70, line width=0.4pt] (sale) -- (entra)
    node[midway, above=7pt, sloped, etiqueta] {$L$};
  \node[etiqueta, anchor=west] at (3.2,3.55) {$X(\alpha,\eta)=\displaystyle\int_{0}^{L}\rho(l)\,\mathrm{d}l$};

  % Detector
  \fill[tolred!85!black] ($(det)+(-0.07,-0.12)$) rectangle ($(det)+(0.07,0.22)$);
  \node[etiqueta, anchor=north] at ($(det)+(0,-0.15)$) {detector};
  \node[etiqueta, anchor=north west, align=left] at (3.2,3.05) {cada dirección $(\alpha,\eta)$\\ es un píxel de la imagen;\\ $L$: recorrido en la roca};
  \node[font=\small\bfseries] at (-3.2,4.3) {(a)};
\end{scope}

% =========================== (b) Ambigüedad espesor-densidad ===========================
\begin{scope}[xshift=6.55cm, yshift=0.1cm]
  \def\w{0.75}
  % Columna 1: 400 m de roca de 2,65 g/cm3
  \filldraw[fill=tolgrey!85, draw=black!70, line width=0.5pt] (0,0) rectangle (\w,2.4);
  % Columna 2: 530 m de roca de 2,00 g/cm3
  \filldraw[fill=tolgrey!35, draw=black!70, line width=0.5pt] (1.8,0) rectangle ({1.8+\w},3.18);
  % Columna 3: 450 m con una cavidad de 50 m y 400 m de roca de 2,65 g/cm3
  \filldraw[fill=tolgrey!85, draw=black!70, line width=0.5pt] (3.6,0) rectangle ({3.6+\w},2.7);
  \filldraw[fill=white, draw=tolblue!80!black, line width=0.5pt] (3.6,1.2) rectangle ({3.6+\w},1.5);

  % Muones incidentes y transmitidos
  \foreach \x/\top in {0/2.4, 1.8/3.18, 3.6/2.7} {
    \draw[->, tolpurple, line width=0.6pt] ({\x+0.5*\w},{\top+0.65}) -- ({\x+0.5*\w},{\top+0.08});
    \draw[->, tolpurple, line width=0.6pt] ({\x+0.5*\w},-0.08) -- ({\x+0.5*\w},-0.5);
    \fill[tolred!85!black] ({\x+0.5*\w-0.22},-0.62) rectangle ({\x+0.5*\w+0.22},-0.54);
  }

  % Rótulos de cada columna
  \node[etiqueta, anchor=north] at ({0.5*\w},-0.7) {$400$~m\\ $2{,}65$~g\,cm$^{-3}$};
  \node[etiqueta, anchor=north] at ({1.8+0.5*\w},-0.7) {$530$~m\\ $2{,}00$~g\,cm$^{-3}$};
  \node[etiqueta, anchor=north] at ({3.6+0.5*\w},-0.7) {$450$~m,\\ cavidad\\ de $50$~m};
  \node[etiqueta, text=tolblue!80!black, anchor=west] at ({3.6+\w+0.08},1.35) {cavidad};

  % Opacidad común
  \draw[black!60, line width=0.3pt] (-0.1,-1.9) -- (4.45,-1.9);
  \node[etiqueta] at (2.17,-2.25) {misma opacidad, $X=1{,}06\times10^{5}$~g\,cm$^{-2}$:\\ mismo conteo esperado en el detector};
  \node[etiqueta, text=tolpurple] at (2.17,4.2) {flujo incidente};
  \node[font=\small\bfseries] at (-0.35,4.2) {(b)};
\end{scope}
\end{tikzpicture}
\end{document}
